\documentclass[aspectratio=169,10pt,spanish]{beamer}

\input{../../../_preambulo_beamer.tex}
\input{../../../_comandos_beamer.tex}

\selectlanguage{spanish}

\title{MA1001B - Modelación Estadística para la Toma de Decisiones}
\subtitle{Lección 17: Modelo de Poisson para eventos raros}
\author{}
\institute{Tecnol\'ogico de Monterrey}
\date{}

\begin{document}

\begin{frame}[plain]
  \vspace{1.2cm}
  {\Large\bfseries MA1001B - Modelación Estadística para la Toma de Decisiones\par}
  \vspace{0.7cm}
  {\large Lección 17: Modelo de Poisson para eventos raros\par}
  \vspace{0.7cm}
  {\color{TecRojo}\rule{\textwidth}{1.2pt}}
  \vspace{0.8cm}

  Facultad MA1001B\\[0.3cm]
  Periodo\\[0.3cm]
  Tecnol\'ogico de Monterrey
\end{frame}

\begin{frame}{Conteos en un turno}
  \begin{alertblock}{Pregunta guía}
    Si las paradas ocurren a tasa media $\lambda=3.2$ por turno, ¿cuáles son
    $P(X=0)$ y $P(X\ge 6)$, y qué ocurre si los eventos se agrupan?
  \end{alertblock}

  \vspace{0.6em}
  Al final de esta lección, usted puede:
  \begin{itemize}
    \item usar la propiedad Poisson de media igual a varianza;
    \item calcular $P(X=0)$ y $P(X\ge 6)$;
    \item comparar esos valores con una simulación de semilla $42$;
    \item explicar la sobredispersión cuando los eventos se agrupan.
  \end{itemize}

  \vfill
  \scriptsize Todos los turnos usados hoy son sintéticos. No se usan datos reales de empresa.
\end{frame}

\begin{frame}{Poisson($\lambda=3.2$)}
  \small
  $X$ es el número de eventos en un turno.
  \[
    E[X]=\lambda=3.2,\qquad \Var(X)=\lambda=3.2,
  \]
  \[
    P(X=0)=e^{-3.2}=0.040762,\qquad P(X\ge 6)=0.105408.
  \]
\end{frame}

\begin{frame}[fragile]{Paso 1: FMP de Poisson}
  \begin{columns}[T]
    \begin{column}{0.42\textwidth}
      \small
      \begin{lstlisting}
model = stats.poisson(mu=3.2)
p0 = model.pmf(0)
p6 = model.sf(5)
      \end{lstlisting}

      \begin{block}{Salida verificada}
        $\lambda=3.200000$\\
        $P(X=0)=0.040762$\\
        $P(X\ge 6)=0.105408$
      \end{block}
    \end{column}
    \begin{column}{0.58\textwidth}
      \centering
      \includegraphics[width=\textwidth]{../figuras/poisson_01_fmp.png}
      \vspace{0.2em}
      \scriptsize FMP de Poisson del Paso 1
    \end{column}
  \end{columns}
\end{frame}

\begin{frame}{Diez mil turnos simulados}
  \small
  Semilla $42$, $n=10{,}000$ extracciones de Poisson:
  \[
    \widehat{P}(X=0)=0.040000,\qquad
    \widehat{P}(X\ge 6)=0.109000.
  \]
  La media simulada $3.209300$ y la varianza $3.256094$ se sitúan cerca de
  $\lambda=3.2$.
\end{frame}

\begin{frame}[fragile]{Paso 2: Comprobación por simulación}
  \begin{columns}[T]
    \begin{column}{0.40\textwidth}
      \small
      \begin{lstlisting}
rng = np.random.default_rng(42)
x = rng.poisson(3.2, 10000)
      \end{lstlisting}

      \begin{block}{Salida verificada}
        $P(X=0)$ simulada $=0.040000$\\
        $P(X\ge 6)$ simulada $=0.109000$\\
        Media $=3.209300$
      \end{block}
    \end{column}
    \begin{column}{0.60\textwidth}
      \centering
      \includegraphics[width=\textwidth]{../figuras/poisson_02_simular_turnos.png}
      \vspace{0.2em}
      \scriptsize Histograma frente a FMP del Paso 2
    \end{column}
  \end{columns}
\end{frame}

\begin{frame}{Paso 3: El agrupamiento produce sobredispersión}
  \begin{columns}[T]
    \begin{column}{0.42\textwidth}
      \small
      Turnos en calma $\lambda=2.0$ ($70\%$) y turnos ocupados $\lambda=6.0$
      ($30\%$) conservan la media $3.2$ pero elevan la varianza a $6.560000$
      en teoría y $6.811382$ en la mezcla con semilla $42$.

      \vspace{0.4em}
      \textbf{Límite:} si los eventos se agrupan, la varianza excede
      $\lambda$. La Lección 18 modela tiempos de espera hasta el primer o el
      $r$-ésimo éxito.
    \end{column}
    \begin{column}{0.58\textwidth}
      \centering
      \includegraphics[width=\textwidth]{../figuras/poisson_03_sobdispersion.png}
      \vspace{0.2em}
      \scriptsize Sobredispersión del Paso 3
    \end{column}
  \end{columns}
\end{frame}

\begin{frame}{Verificación de conceptos}
  \begin{enumerate}
    \item ¿Cuál es $\mathrm{Var}(X)$ para una variable Poisson($3.2$)?
    \item Calcule $P(X=0)$ a partir de $e^{-3.2}$.
    \item ¿Por qué una tasa simulada de $0.040000$ no sustituye a $0.040762$?
    \item ¿Por qué una mezcla de $\lambda=2$ y $\lambda=6$ sobredispersa?
  \end{enumerate}

  \vspace{0.6em}
  \begin{block}{Conexión con la Lección 18}
    La sesión puede continuar directamente con tiempos de espera geométricos
    y binomiales negativos.
  \end{block}

  \vfill
  \scriptsize Fuente: OpenStax, \textit{Introductory Business Statistics 2e},
  Sección 4.4. CC BY-NC-SA.
\end{frame}

\appendix



\end{document}
